(1) 在某一瞬时火箭静止系内，发射尾流静质量$dm_e>0$、速率$u$，其能量和动量大小为$\gamma_u dm_ec^2$与$\gamma_udm_eu$，$\gamma_u=(1-u^2/c^2)^{-1/2}$。发射后火箭质量为$M+dM$（$dM<0$），微小前向速度$dv$。能量守恒到一阶给$-dM=\gamma_udm_e$，动量守恒给$M\,dv=\gamma_udm_eu$，故$dv=-u\,dM/M$；这里$dv$专指共动系的微小相对速度，而非固定实验室速度差。
(2) 用快速度$\eta=\operatorname{artanh}(v/c)$处理速度复合：局域增量$d\eta=dv/c=-(u/c)dM/M$。从$M_0$积分至$M_f$，若火箭初始静止则$\eta_f=(u/c)\ln(M_0/M_f)$，最终$v_f=c\tanh[(u/c)\ln(M_0/M_f)]$。以快速度而非线性速度相加，保证全部质量比下$v_f<c$。
(3) 若$|\eta_f|\ll1$，$\tanh\eta_f\simeq\eta_f$，回复经典齐奥尔科夫斯基公式$v_f\simeq u\ln(M_0/M_f)$。对极限$u\uparrow c$，有$v_f/c=\tanh\ln(M_0/M_f)=(M_0^2-M_f^2)/(M_0^2+M_f^2)$，即理想光子火箭方程；有限质量比始终不能达到光速。
(4) 定义目标快速度$\eta_f=\operatorname{artanh}(v_f/c)$，反求$M_0/M_f=\exp[(c/u)\operatorname{artanh}(v_f/c)]$。若质量比趋一，指数趋零，所以$v_f\to0$；若$u$固定且$v_f\to c^-$，质量比发散。该模型将尾流能量包含在火箭质量损失内，因此不能再无端把排气静质量误认为$-dM$而遗漏$\gamma_u$。
